%
% Die Stammfunktion von 1/theta
%
\subsection{Die Stammfunktion von $1/\vartheta$
\label{buch:intalg:wurzel:subsection:stammfunktiontheta}}
Durch Nachrechnen der Ableitung von
\[
\Theta
=
\frac{1}{\!\sqrt{a}}
\log\Bigl(
2ax + b - 2\!\sqrt{a} \vartheta
\Bigr)
\]
bekommt man
\begin{align*}
\Theta'
&=
-\frac{1}{\!\sqrt{a}}
\cdot
\frac{1}{2ax+b-2\!\sqrt{a}\vartheta}
\cdot
(2a - 2\!\sqrt{a}\vartheta')
\\
&=
-\frac{1}{\!\sqrt{a}}
\cdot
\frac{1}{2ax+b-2\!\sqrt{a}\vartheta}
\cdot
\biggl(
2a -2\!\sqrt{a}\frac{2ax+b}{2\vartheta}
\biggr)
\\
&=
-\frac{1}{\!\sqrt{a}}
\cdot
\frac{1}{2ax+b-2\!\sqrt{a}\vartheta}
\cdot
\frac{
4a\vartheta -2\!\sqrt{a}(2ax+b)
}{
2\vartheta
}
\\
&=
-\frac{1}{\!\sqrt{a}}
\cdot
\frac{1}{2ax+b-2\!\sqrt{a}\vartheta}
\cdot
\frac{
2\!\sqrt{a}\cdot
\bigl(
\!\sqrt{a}\vartheta
-
2ax-b
\bigr)
}{
2\vartheta
}
=
\frac{1}{\vartheta},
\end{align*}
Die Funktion $\Theta$ ist also eine Stammfunktion von $1/\vartheta$.
Sie liegt in einem Erweiterungskörper, dem die Konstante $\!\sqrt{a}$
und der Logarithmus $\log(2ax+b-2\!\sqrt{a}\vartheta)$ hinzugefügt
werden ist.
Damit der Logarithmus hinzugefügt werden kann, muss auch noch $b$
hinzugefügt werden.
Somit ist
\[
\Theta
\in
\mathbb{Q}\bigl(x,\theta,\!\sqrt{a},b,\log(\!\sqrt{a}\vartheta-2ax-b)\bigr)
\]
eine elementare Stammfunktion von $1/\vartheta$.

